\section{Introduction}
\label{sec:intro}

A tool call is an executive act. When an LLM agent reads a file, writes to a database, sends an email, executes a shell command, or pushes a commit, it takes an action on behalf of a user, against an external system, with consequences in the world. The agent receives delegated authority through a system prompt or developer credential; it dispatches that authority against a resource it does not own; it produces an effect that, if wrong, must be undone.

The agentic-AI safety literature has converged on a particular framing. Safety is treated as a property of the model: how to make the model want the right thing, refuse the wrong thing, and behave consistently between evaluation and deployment. Constitutional AI~\cite{baiConstitutionalAI2022} and RLHF~\cite{ouyangRLHF2022} shape preferences at training time. Scalable oversight~\cite{leikeScalableOversight2018,bowmanScalableOversight2022} and debate~\cite{christianoDebate2018} extend supervision past the operator's direct capability. Alignment-faking~\cite{hubingerAlignmentFaking2024} and situational-awareness~\cite{berglundSituationalAwareness2023} work sharpen the threat model. Each thread treats the model as the object of safety analysis.

We argue that safety is also, separately, a property of the substrate that admits the tool call. The substrate may decline to execute a reversible call without typed rollback evidence or a destructive call without bilateral consent, regardless of whether training-layer alignment has fully succeeded and regardless of whether the model has evaded its training signal. Admission-layer safety composes with training-layer safety rather than displacing it (Figure~\ref{fig:orthogonal-composition}): a model whose preferences are well-shaped by training will, in the typical case, issue admissible envelopes, and the substrate's role is to bound the operational consequence in the cases where training has not held. The grammar of admission is independent of the grammar of generation.

The structural argument is short, but conditional. A tool call has the shape of a delegated executive act: a typed envelope, an authorisation chain, a target, an effect, and (in the destructive case) a need for reversal or bilateral consent. A substrate that mediates every call through a correctly classified single-envelope path can require a positive time-to-live (TTL) witness and a typed rollback receipt slot for reversible effects, and destructive-class bilateral consent for effects without a credible inverse. The model may be misaligned; in the mediated path it does not choose the action class, mutate the registry, bypass the hook, or substitute its own verifier for the substrate's.

\paragraph{Contributions.}
This paper makes three contributions, all structural.
\begin{itemize}
\item We frame the agent tool call as a five-part executive act (authorisation chain, subject, effect, lifetime, reversal expectation), and we show that the same five components appear in adjacent formally analysed domains (endpoint security response, temporally bounded administrative orders).
\item We name a substrate-layer safety property orthogonal to the training-layer, evaluation-layer, and interpretability-layer stack: typed admission with TTL by construction, rollback receipt slot for reversible classes, and bilateral cosignature on destructive classes. We give the grammar and cite the formal substrate that carries the broader admission and refinement machinery~\cite{programmableSovereignty2026}.
\item We identify the threat model boundary for the substrate-layer guarantee: mediated, correctly classified, single-envelope execution under independent cosigners, an intact registry, no bypass path, and working rollback executors. We separately name the cases outside that boundary: mis-classification, registry compromise, multi-agent collusion, host-side willing signatures, time-of-check/time-of-use races, admission-hook denial of service, and faulty rollback code.
\end{itemize}

The contribution is structural rather than empirical. The paper does not introduce a new training-time technique or a benchmark. It proposes that two safety layers exist, that they compose, and that the runtime-substrate layer can enforce a bounded admission property under adversarial assumptions about the model that training-time interventions cannot match by themselves.

The argument proceeds as follows. Section~\ref{sec:background} surveys the existing stack and identifies the gap. Section~\ref{sec:executive-acts} develops the executive-act framing. Section~\ref{sec:grammar} states the admission grammar. Section~\ref{sec:impl} gives a worked example. Section~\ref{sec:threat} states the bounded threat model. Section~\ref{sec:discussion} situates the construction next to existing alignment-research threads. Section~\ref{sec:limitations} names the bounds, and Section~\ref{sec:conclusion} concludes.
